\section{The composite kernels on the families}

Throughout this section $\mathfrak g=\mathfrak{su}(N)$ at leading order in $N$ and $O$ is a single trace, with the conventions of Section 10.7. The projection to cohomology on the two-field part is the evaluation on the free equations of motion, $|p|^2=0$ on the scalar legs, the Weyl equation on the gaugino legs and transverse on-shell gauge legs, which removes the $s$-exact terms $s(\text{antifield}\times\text{field})$; total derivatives are kept, as descendants through $\omega$. The letters are written in the components of Definition \ref{def:c8-twocomp}.

\begin{definition}[Test families]
\label{def:c11-families}
Let $n_m:=\sigma^m_{1\dot1}=(0,0,1,i)$, a null vector, and $D_n:=n_m(\partial_m+[A_m,-])$, the covariant derivative $\iota_nd_{A|_V}$ along $n$; let $Z=\iota_wA$ with $w\in W_{\mathbb R}^\perp\otimes\mathbb C$, $(w,w)=0$; let $\psi:=\lambda_1{}^i$ for a fixed $i$, the component $\alpha=1$ of $\lambda|_{S_+\otimes W}$; and let $F_{11}$ be the component $\alpha=\beta=1$ of $(\sigma^m\bar\sigma^n)_\alpha{}^\gamma\varepsilon_{\gamma\beta}F_{mn}$, $F_{mn}=\partial_mA_n-\partial_nA_m+[A_m,A_n]$, whose two-field part is a component of $dA|_V$ and whose three-field part is a component of $[A|_V\wedge A|_V]$. The test families are the spans of the following single traces and their gauge completions:
\begin{itemize}
\item the scalars: $\operatorname{tr}(\iota_{w_1}A\cdots\iota_{w_L}A)$, the bilinears, and the $\mathfrak{su}(2)$ sector of Proposition \ref{prop:c10-su2} for $L\le8$;
\item twist two: $\operatorname{tr}\big(D_n^kZ\,D_n^{S'-k}Z\big)$, $0\le k\le S'$ (Proposition \ref{prop:c10-twist});
\item the $\mathfrak{su}(1|1)$ sector: single traces of the letters $Z$ and $\psi$ of Proposition \ref{prop:c10-su11}, $L\le6$;
\item the fermion $\mathfrak{sl}(2)$ sector: $\operatorname{tr}\big(D_n^k\psi\,D_n^{S'-k}\psi\big)$;
\item two gaugino components: $\operatorname{tr}\big(D_n^k\psi^{(3)}\,D_n^{S'-k}\psi^{(4)}\big)$ with $\psi^{(3)}=\lambda_1{}^3$, $\psi^{(4)}=\lambda_1{}^4$, together with $\operatorname{tr}\big(D_n^kZ'\,D_n^{S'-k}F_{11}\big)$, $Z'$ the scalar letter with the R-symmetry weights of $\psi^{(3)}\psi^{(4)}$, with which the family closes;
\item the gluon $\mathfrak{sl}(2)$ sector: $\operatorname{tr}\big(D_n^kF_{11}\,D_n^{S'-k}F_{11}\big)$;
\item the mixed gluon--gaugino family: $\operatorname{tr}\big(D_n^kF_{11}\,D_n^{S'-k}\psi\big)$.
\end{itemize}
Twist two, the fermion and the gluon $\mathfrak{sl}(2)$ sectors are families of Proposition \ref{prop:c11-theta-sector} for the sectors of $A|_{W_{\mathbb R}^\perp}$, of $\lambda$ and of $A|_V$.
\end{definition}

The two-field part of a gauge-invariant operator with covariant derivatives receives the pairs of its two leading letters with two outputs and the pairs of a letter $A|_V$ from a covariant derivative, or from the three-field part of a field strength, with another letter, with one output (Proposition \ref{prop:c10-twist}). For the composite pair kernels these are the leg sets $A|_{W_{\mathbb R}^\perp}$ of $(A|_V,A|_{W_{\mathbb R}^\perp})$, $\lambda$ of $(A|_V,\lambda)$ and $A|_V$ of $(A|_V,A|_V)$ of Table \ref{tab:c11-fieldpairs}, all with $k=1$, with derivatives.

\begin{proposition}[Spectra from the harmonic action]
\label{prop:c11-targets}
On the two-site graded cyclic states, single traces of two letters with $S'$ derivatives $D_n$ in total, the eigenvalues of $\sum_lH_{l,l+1}$ of [B03c, (1.1)--(1.3)], computed from the harmonic action [B03c, (4.27)--(4.29)] on the oscillator letters [B03c, (4.9)], are:
\begin{itemize}
\item (i) scalars $D_n^kZ$: $4h(S)$, $S\le S'$ even [B03c, (2.31)];
\item (ii) gluons $D_n^kF_{11}$: $4h(S+2)$, $S\le S'$ even, so that $\operatorname{tr}(F_{11}F_{11})$ has the Konishi value $6$;
\item (iii) one gaugino component $D_n^k\psi$: $4h(S+1)$, $S\le S'$ odd;
\item (iv) two gaugino components $\psi^{(3)}$, $\psi^{(4)}$ with $Z'F_{11}$: $0,6$ at $S'=0$; $0,6,6,6$ at $S'=1$; $0,6,6,6,6,\tfrac{25}3$ at $S'=2$; $0$, four times $6$ and three times $\tfrac{25}3$ at $S'=3$; $0$, four times $6$, four times $\tfrac{25}3$ and $\tfrac{49}5$ at $S'=4$;
\item (v) $(F_{11},\psi)$: $6$; $6,6$; $6,6,\tfrac{25}3$; $6,6,\tfrac{25}3,\tfrac{25}3$; $6,6,\tfrac{25}3,\tfrac{25}3,\tfrac{49}5$ at $S'=0,\dots,4$.
\end{itemize}
On the $\mathfrak{su}(1|1)$ sector the two-site Hamiltonian is one minus the graded exchange (Proposition \ref{prop:c10-su11}), with the cyclic spectra of Table \ref{tab:c11-su11}.
\end{proposition}

\begin{proof}
Exact evaluation of [B03c, (4.27)--(4.29)], with the harmonic numbers of the coefficients $c_{n,0,0}=h(\tfrac n2)$, on bases of the families closed under $H_{12}$, with the graded exchange for $H_{21}$ and the projection to graded cyclic states, and exact diagonalisation. The same evaluation gives [B03c, (2.31)] on twist two and the values $0$, twenty times, and $6$ on the scalar bilinears, which fixes the central charge per site, the signs of the fermionic oscillators and the graded cyclicity.
\end{proof}

\begin{tabular}{l|l}
$(L,M)$ & eigenvalues on the graded cyclic states of $\operatorname{tr}(Z^{L-M}\psi^M)$\\
\hline
$M\le1$ & $0$\\
$(3,2)$, $(3,3)$ & $6$\\
$(4,2)$, $(4,3)$ & $4$\\
$(5,2)$, $(5,3)$ & $5\pm\sqrt5$\\
$(5,4)$, $(5,5)$ & $10$\\
$(6,2)$ & $2$, $6$\\
$(6,3)$ & $2$, $6$, $8$, $8$\\
$(6,4)$ & $8$, $8$, $8$\\
$(6,5)$ & $8$\\
\end{tabular}
\tabcaption{The eigenvalues of $\sum_lH_{l,l+1}$ on the $\mathfrak{su}(1|1)$ sector, $L\le6$; the sectors $(2,2)$, $(4,4)$, $(6,6)$ have no graded cyclic state.\label{tab:c11-su11}}

\begin{theorem}[The composite pair kernels on the families]
\label{thm:c11-families}
\begin{itemize}
\item (i) (The two-site density.) For every ordered adjacent pair $(X_a,X_b)$ of letters among $A|_{W_{\mathbb R}^\perp}$, $\lambda|_{S_+\otimes W}$, $\lambda|_{S_-\otimes W^*}$ without derivatives, on the free shell, and for two letters $A|_{W_{\mathbb R}^\perp}$ also off the shell and for all their components, the planar composite pair kernels have, on the same word $X_aX_b$, the translation-invariant part $-\tfrac43$ units times $X_aX_b$ along $\Delta$ and the value $-\tfrac16$ units times $\partial_b(X_aX_b)$ along $P_b$; on the exchanged word and on the trace they vanish.
\item (ii) On the two-field parts of the families of Definition \ref{def:c11-families}, on the free shell, the sum over the pairs of the translation-invariant parts of the composite pair kernels along $\Delta$ is the multiple of the identity of Table \ref{tab:c11-results}, which is $-\tfrac23$ units times the letter count, and the sum of the composite pair kernels along $n$ is $-\tfrac13$ units times $\partial_n$; on twist two for $S'\le4$ the sum along every $b\in V$ is $-\tfrac13$ units times $\partial_b$.
\item (iii) Hence (b) along $n$ and (c) of Definition \ref{def:c11-P7} hold on each family in the ranges of Table \ref{tab:c11-results}, and (b) along every $b$ holds on twist two for $S'\le4$.
\end{itemize}
\end{theorem}

\begin{proof}
Exact evaluation of the special-edge composite cycles by Theorem \ref{thm:c11-simplex}, Propositions \ref{prop:c11-genfun} and \ref{prop:c11-k1}, with the colour words reduced to the planar adjacent form of Lemma \ref{lem:c10-colour} (iv), on the leg sets of Table \ref{tab:c11-fieldpairs} with the derivatives of the families, the outputs evaluated on the free shell (Remark \ref{rem:c11-method}). On the families with derivatives the outputs depend on the output momenta through $(n,p)$, every derivative being contracted with $n$ and $(n,n)=0$; the sums are diagonal in the bases of the families. (iii) is (i), (ii) and the values of the letter count: $2L$ on $L$ letters without derivatives, $4+S'$ on twist two and on two gaugino letters, $5+S'$ on $Z'F_{11}$ and on $(F_{11},\psi)$, $6+S'$ on two gluon letters, a letter $F_{11}$ counting $3$ through its two-field part $dA|_V$.
\end{proof}

\begin{tabular}{l|l|l|l}
family & range & sum of the translation-invariant parts along $\Delta$ & spectrum of $D_1$ at $\theta=\theta_\lambda=\theta_V=-\tfrac13$, in units\\
\hline
scalars & bilinears; $\mathfrak{su}(2)$, $L\le8$ & $-\tfrac43L$ & $0$, $-48$; $-8$ times the Heisenberg spectra\\
twist two & $S'\le6$ & $-\tfrac23(4+S')$ & $0$, $-48$, $-\tfrac{200}3$, $-\tfrac{392}5$: $-32h(S)$, $S\le S'$ even\\
$\mathfrak{su}(1|1)$ & $L\le6$ & $-\tfrac43L$ & $-8$ times the $\mathfrak{su}(1|1)$ spectra\\
fermion $\mathfrak{sl}(2)$ & $S'\le5$ & $-\tfrac23(4+S')$ & $-48$, $-\tfrac{200}3$, $-\tfrac{392}5$: $-32h(S+1)$, $S\le S'$ odd\\
$\psi^{(3)}\psi^{(4)}$ with $Z'F_{11}$ & $S'\le2$ & $-\tfrac23(4+S')$ on $\psi^{(3)}\psi^{(4)}$, $-\tfrac23(5+S')$ on $Z'F_{11}$ & $0,-48$; $0,-48,-48,-48$; $0$, four times $-48$, $-\tfrac{200}3$\\
gluon $\mathfrak{sl}(2)$ & $S'\le2$ & $-\tfrac23(6+S')$ & $-48$, $-\tfrac{200}3$: $-32h(S+2)$, $S\le S'$ even\\
$(F_{11},\psi)$ & $S'\le4$ & $-\tfrac23(5+S')$ & at $S'=4$: $-48$, $-48$, $-\tfrac{200}3$, $-\tfrac{200}3$, $-\tfrac{392}5$\\
\end{tabular}
\tabcaption{The composite pair kernels and the spectrum of $D_1$ on the families: the third column is identity (c), the fourth is $-8$ units times the eigenvalues of Proposition \ref{prop:c11-targets}, of Proposition \ref{prop:c10-su2} on the $\mathfrak{su}(2)$ sector and of Table \ref{tab:c11-su11} on the $\mathfrak{su}(1|1)$ sector.\label{tab:c11-results}}

\begin{corollary}[$D_1$ and $P_1$ on the families]
\label{cor:c11-D1fam}
Under the hypotheses of Definition \ref{def:c10-hyp}, on the cohomology of the two-field part of every family of Definition \ref{def:c11-families}, with every free direction of Definition \ref{def:c11-free},
\[
D_1\cong4\,(\text{the map})-\sum_a\big(\theta_{X_a}+\tfrac13\big)\,\text{units}\cdot v.X_a,\qquad P_1\cong-\sum_a\big(\theta_{P,X_a}+\tfrac13\big)\,\text{units}\cdot\partial_bX_a,
\]
with the map of Theorem \ref{thm:c10-brst}, where $\theta_{X_a}$ is $\theta$, $\theta_\lambda$ or $\theta_V$ and $\theta_{P,X_a}$ is $\theta_P$, $\theta_{P,\lambda}$ or $\theta_{P,V}$ according to the sector of the letter $X_a$, $v.X_a$ and $\partial_bX_a$ act on the letter $X_a$ alone, and the directions $Q\int(\iota_xA|_V,c^*)$, $Q\int(\rho(x)\lambda^*,\lambda^*)$ contribute $s$-exact terms. Under the hypotheses of Theorem \ref{thm:c11-arity2}, twist two gives $\theta=-\tfrac13$, the fermion $\mathfrak{sl}(2)$ sector $\theta_\lambda=-\tfrac13$ and the gluon $\mathfrak{sl}(2)$ sector $\theta_V=-\tfrac13$ (Proposition \ref{prop:c11-theta-sector}); then $D_1$ induces four times the map of Theorem \ref{thm:c10-brst} on every family, with the spectra of Table \ref{tab:c11-results}, $-8$ units times the eigenvalues of $\sum_lH_{l,l+1}$, which is twice the one-loop anomalous dimensions; and $P_1\cong0$ in the translation-covariant scheme, while for equal coefficients $\theta_P$ along $P_b$, $P_1\cong(\theta-\theta_P)$ units times $\partial_b$ with $\theta=-\tfrac13$.
\end{corollary}

\begin{proof}
Corollary \ref{cor:c11-D1local} and Theorem \ref{thm:c11-families}: on the two-field part $\hbar F_1^{(1)}(\Delta)$ acts through its quadratic part, the bracket of the sector part of the sector of $X_a$ with $X_a$ being $v.X_a$ and the others vanishing. As in the proof of Theorem \ref{thm:c11-P7} (i), letter by letter, the position part of the composite pair kernels and of the bracket is $-2(\tfrac13+\theta_{X_a})\partial_b$ units and the translation-invariant part is $-2(\tfrac13+\theta_{X_a})w_{X_a}$ units; along $P_b$ the composite pair kernels give $-\tfrac13$ units times $\partial_b$ and the bracket with $\hbar F_1^{(1)}(P_b)$ gives $-\theta_{P,X_a}$ units times $\partial_bX_a$. The $\kappa$-directions are $s$-exact (Proposition \ref{prop:c11-arity1} (ii)). The rest is Proposition \ref{prop:c11-theta-sector}, Theorem \ref{thm:c11-P7} and Corollary \ref{cor:c11-theta} (iii), (iv).
\end{proof}

\begin{remark}[Method of computation]
\label{rem:c11-method}
Each special-edge composite cycle is evaluated exactly: the Gaussian loop integral and the integrals of \eqref{eq:c11-simplex} over the ordered subsimplices of the region where $\alpha_\sigma$ is the smallest coordinate, the letter derivatives through Proposition \ref{prop:c11-genfun}, the position through the injected momentum $r$ (Proposition \ref{prop:c11-position}), the chains with $k=1$ through Proposition \ref{prop:c11-k1}, the fibre algebra in explicit realisations of $V_{10}$ and $\Sigma$ over $\mathbb Q(i)$, and the colour words by Lemma \ref{lem:c10-colour}, with its planar adjacent rule also checked in realisations of $\mathfrak{su}(2)$, $\mathfrak{su}(3)$ and $\mathfrak{gl}(3)$. In this evaluation the factor $(-1)^{[c<\sigma]}$ of Definition \ref{def:c11-sp} is part of the chain convention when the special edge ends at $X_b$ and is applied as the factor $-1$ when the special edge joins two vertices of $I$. A second evaluation works in momentum space with the kernel form of Lemma \ref{lem:c11-kernels} (ii), $v.$ acting on the Fourier components as an Euler operator plus constants, Feynman parameters for the edges $P$ and the special Schwinger time integrated last, and applies the sign of Definition \ref{def:c11-sp} as written. The two evaluations give the same off-shell entries with at most two outputs and no letter derivatives, for the pairs of bosonic letters and for the pairs with a gaugino letter or output, whose values lie in $\mathbb Q+\mathbb Q\log2+\mathbb Q\log3$ for $k\le2$, and the same family matrices on the scalars, on twist two for $S'\le4$, on the $\mathfrak{su}(1|1)$ sector, on the fermion $\mathfrak{sl}(2)$ sector for $S'\le5$ and on the gluon $\mathfrak{sl}(2)$ sector for $S'\le2$; the first evaluation extends twist two to $S'\le6$, and the second evaluation gives the two families with $F_{11}$ and a gaugino letter. A third evaluation, a Taylor expansion in momentum space at generic null leg momenta, gives the sums on twist two for $S'\le4$ along a generic $b$. Every integral is a four-dimensional Schwinger integral with the cutoff $\varepsilon$ on every edge.
\end{remark}

